\documentclass[11pt]{article}
\usepackage[margin=1in]{geometry}
\usepackage{lmodern}
\usepackage{hyperref}
\usepackage{enumitem}
\usepackage{amsmath,amssymb,amsthm,mathtools}
\usepackage[T1]{fontenc}
\usepackage[utf8]{inputenc}
\title{\textbf{Theory-mode report: Quantitative Estimates for Mean-Field Limits and Correlation Functions through a Duality Framework}}
\author{Generated from Scholar Inbox digest}
\date{Digest date: 2026-05-05}
\begin{document}
\maketitle
\section*{Paper information}
\begin{itemize}[leftmargin=1.5em]
\item \textbf{Title:} Quantitative Estimates for Mean-Field Limits and Correlation Functions through a Duality Framework
\item \textbf{Authors:} Nadia Khoury, P. -E. Jabin
\item \textbf{arXiv:} \href{https://arxiv.org/abs/2605.02058}{2605.02058}
\item \textbf{Scholar digest score:} 80.0
\end{itemize}
\section*{Executive summary}
This report summarizes the highest-scoring paper in the Scholar Inbox digest dated \textbf{2026-05-05}. The abstract states the core claim as follows: \emph{We investigate the mean-field limit for interacting particle systems through a duality-based framework and obtain quantitative estimates on the convergence of marginals as well as on correlation functions. In particular, for merely square-integrable interaction forces, we derive the natural fluctuation-scale rate \$\textbackslash\{\}mathcal\{O\}(N\textasciicircum{}\{-1/2\})\$. By introducing an iterative argument on the hierarchy of dual cumulants, we leverage this bound to recover the optimal mean-field rate \$\textbackslash\{\}mathcal\{O\}(N\textasciicircum{}\{-1\})\$ and to obtain robust estimates on the dual cumulants, at the expense of corresponding regularity assumptions on the interaction kernel. Finally, using the relation between dual and direct correlations, we transfer these bounds to direct cumulants, yielding refined information on correlations and deviations from chaos.}.

\section*{Problem setup and intuition}
The introduction reinforces this lens: the paper appears motivated by a gap between practical behavior and theoretical understanding.

\section*{Proof architecture}
The technical core appears to build the surrogate quantity and the chain of lemmas needed to propagate local control into the final theorem: the paper follows a standard theory-paper structure with structural reduction plus quantitative bounds.

\section*{Takeaway}
The key reusable lesson is to define the right object, find the structural reduction, and quantify the error terms clearly.
\end{document}
